\documentclass[11pt]{article}
\usepackage[margin=1in]{geometry}
\usepackage{amsmath, amssymb}
\usepackage{enumitem}
\usepackage{parskip}

\title{ORF 309 Homework 3 -- Question 4: Love Triangles}
\author{}
\date{}

\begin{document}
\maketitle

\section*{Problem Setup}

A dating website has $n$ people signed up. For every pair of people, the website independently flips a coin that comes up heads with probability $p$. If heads, the pair is matched. A \textbf{love triangle} is a set of three people who are all mutually matched.

\section*{Part (a): Expected Number of Hookups for a Given Individual}

\textbf{Step 1: Identify the random experiment.}

Fix a particular person, call them Person 1. There are $n - 1$ other people on the site. For each of these $n-1$ people, the website independently flips a coin with probability $p$ of heads.

\textbf{Step 2: Define indicator variables.}

For each person $j \in \{2, 3, \ldots, n\}$, define:
\[
X_j = \begin{cases} 1 & \text{if Person 1 is matched with Person } j \\ 0 & \text{otherwise} \end{cases}
\]
Each $X_j$ is a Bernoulli random variable with $P(X_j = 1) = p$, so $E[X_j] = p$.

\textbf{Step 3: Express total hookups as a sum.}

The total number of hookups Person 1 has is:
\[
H = \sum_{j=2}^{n} X_j
\]

\textbf{Step 4: Apply linearity of expectation.}
\[
E[H] = E\left[\sum_{j=2}^{n} X_j\right] = \sum_{j=2}^{n} E[X_j] = \sum_{j=2}^{n} p = (n-1)p
\]

\begin{center}
\framebox{\textbf{Answer: } $E[\text{hookups}] = (n-1)p$}
\end{center}

\section*{Part (b): Probability of More Than One Hookup}

\textbf{Step 1: Identify the distribution.}

From Part (a), $H = \sum_{j=2}^{n} X_j$ where each $X_j$ is an independent Bernoulli($p$). The sum of $n-1$ independent Bernoulli($p$) variables follows a \textbf{Binomial distribution}:
\[
H \sim \text{Binomial}(n-1,\; p)
\]

\textbf{Step 2: Use the complement.}
\[
P(H > 1) = 1 - P(H = 0) - P(H = 1)
\]

\textbf{Step 3: Compute $P(H = 0)$.}
\[
P(H = 0) = \binom{n-1}{0} p^0 (1-p)^{n-1} = (1-p)^{n-1}
\]
This is the probability that \emph{none} of the $n-1$ coin flips come up heads.

\textbf{Step 4: Compute $P(H = 1)$.}
\[
P(H = 1) = \binom{n-1}{1} p^1 (1-p)^{n-2} = (n-1)\,p\,(1-p)^{n-2}
\]
This is the probability that exactly one of the $n-1$ flips comes up heads.

\textbf{Step 5: Combine.}
\[
P(H > 1) = 1 - (1-p)^{n-1} - (n-1)\,p\,(1-p)^{n-2}
\]

\begin{center}
\framebox{\textbf{Answer: } $P(\text{more than one hookup}) = 1 - (1-p)^{n-1} - (n-1)\,p\,(1-p)^{n-2}$}
\end{center}

\section*{Part (c): Expected Number of Love Triangles}

\textbf{Step 1: Enumerate all possible triangles.}

A love triangle requires choosing 3 people out of $n$. The number of ways to do this is:
\[
\binom{n}{3} = \frac{n(n-1)(n-2)}{6}
\]

\textbf{Step 2: Define indicator variables for each triple.}

For each triple of people $\{i, j, k\}$, define:
\[
T_{ijk} = \begin{cases} 1 & \text{if } \{i, j, k\} \text{ form a love triangle} \\ 0 & \text{otherwise} \end{cases}
\]
The total number of love triangles on the site is:
\[
L = \sum_{\{i,j,k\}} T_{ijk}
\]
where the sum is over all $\binom{n}{3}$ triples.

\textbf{Step 3: Compute $P(T_{ijk} = 1)$ for a given triple.}

For $\{i, j, k\}$ to be a love triangle, \textbf{all three pairs} must be matched:
\begin{itemize}[nosep]
    \item Pair $(i,j)$ matched with probability $p$
    \item Pair $(i,k)$ matched with probability $p$
    \item Pair $(j,k)$ matched with probability $p$
\end{itemize}
Since the coin flips for different pairs are \textbf{independent}:
\[
P(T_{ijk} = 1) = p \cdot p \cdot p = p^3
\]

\textbf{Step 4: Apply linearity of expectation.}
\[
E[L] = E\left[\sum_{\{i,j,k\}} T_{ijk}\right] = \sum_{\{i,j,k\}} E[T_{ijk}] = \sum_{\{i,j,k\}} p^3 = \binom{n}{3} \cdot p^3
\]

\textit{Note:} Linearity of expectation applies regardless of whether the $T_{ijk}$ are independent---and they are \textbf{not} independent, since triangles can share edges. This is precisely why linearity of expectation is so powerful here.

\textbf{Step 5: Simplify.}
\[
E[L] = \frac{n(n-1)(n-2)}{6} \cdot p^3
\]

\begin{center}
\framebox{\textbf{Answer: } $\displaystyle E[\text{love triangles}] = \frac{n(n-1)(n-2)}{6}\, p^3$}
\end{center}

\end{document}
